\setbeamertemplate{footline}{\hspace{2ex}\scriptsize\textcolor{ink!70}{Historical snapshot: superseded reference/thermals; current validation does not transfer}\hfill\textcolor{ink!60}{\insertframenumber}\hspace{2ex}\vspace{1ex}}
\begin{frame}{Earlier exploratory results: historical context}
  \begin{block}{Superseded reference and thermal treatment}
    The following figures and benchmark summaries retain the earlier calculation
    snapshot. They are not the updated candidate results or the current validation verdict.
  \end{block}
  \vspace{1ex}
  Current results use repaired Fc geometries, uniform solution DFT thermochemistry,
  addressed caches and field-specific validation gates.
\end{frame}

\begin{frame}{Historical redox benchmark (OROP)}
  \begin{columns}[T]
    \begin{column}{0.5\textwidth}
      \textbf{Production-protocol OROP benchmark} ($n=169$):
      \begin{table}\footnotesize
      \begin{tabular}{lcc}
        \toprule
         & ours & OROP impl.-DFT \\
        \midrule
        MAE (V)        & 0.52 & 0.28 \\
        signed bias    & $+0.49$ & --- \\
        Spearman (all) & 0.92 & --- \\
        \bottomrule
      \end{tabular}
      \end{table}
      \vspace{1ex}
      At OROP's own level our organic potentials match theirs ($\sim$0.07 V); the gap sits in
      the Fc reference (4.19 vs 4.66 V). Sourced MeCN set: E1 bias $+0.26$ V, E2 $-0.07$ V.
      \hi{Not calibrated away}.
    \end{column}
    \begin{column}{0.5\textwidth}
      \textbf{Ranking must be judged \emph{within} charge class}
      \begin{table}\footnotesize
      \begin{tabular}{lccc}
        \toprule
        class & $n$ & MAE & Spearman \\
        \midrule
        $+1/0$ couples & 118 & 0.48 & \good{0.79} \\
        $0/-1$ couples & 49 & 0.60 & \good{0.78} \\
        \bottomrule
      \end{tabular}
      \end{table}
      \vspace{1ex}
      \footnotesize Global correlation is inflated because charge classes separate.
      Ranking is assessed within each class; the integrity gate removes category errors.
    \end{column}
  \end{columns}
\end{frame}

% ---------------------------------------------------------------
\begin{frame}{The anion ``failure'' was category errors, not bad physics}
  \begin{columns}[T]
    \begin{column}{0.54\textwidth}
      Some apparent anion failures are category errors: molecules with
      \bad{no intact, bound reduced state}:
      \begin{table}\footnotesize
      \begin{tabular}{llc}
        \toprule
        molecule & exp & RMSD$_\text{neu$\to$an}$ \\
        \midrule
        CCl$_4$    & $-2.57$ & \bad{3.4 \AA} \\
        CH$_2$Br$_2$ & $-3.12$ & 1.3 \AA \\
        (unbound anions) & --- & EA$_\text{gas}<0$ \\
        \bottomrule
      \end{tabular}
      \end{table}
      Dissociative electron attachment (CCl$_4$+e$^-\!\to$CCl$_3^{\bullet}$+Cl$^-$) ---
      scored against irreversible \emph{peak} potentials.
    \end{column}
    \begin{column}{0.46\textwidth}
      \begin{block}{Integrity screen}
        Bound anion (solvated EA $>0$) \emph{and} intact (\hi{bond-graph unchanged}, not
        raw RMSD --- which is confounded by tether flexibility).
      \end{block}
      \vspace{-0.5ex}
      \begin{itemize}\footnotesize
        \item filtering $\rightarrow$ anion Spearman \bad{0.49}$\rightarrow$\good{0.95}
        \item a required necessary-condition gate and stability signal
        \item does \bad{not} test proton transfer, dimerization, or long-term cycling
      \end{itemize}
    \end{column}
  \end{columns}
\end{frame}

% ---------------------------------------------------------------
\begin{frame}{Stability = decomposition free energies}
  \begin{columns}[T]
    \begin{column}{0.54\textwidth}
      \textbf{Disproportionation} $2R \!\rightarrow\! R^{+}\!+R^{-}$ \\ (atom- \&
      electron-conserving; ranks well, $\rho$ 0.80, but \bad{over-estimated by $\sim$23 kJ/mol}):
      \begin{table}\footnotesize
      \begin{tabular}{lr}
        \toprule
        intermediate & $\Delta G_\text{disp}$ (kJ/mol) \\
        \midrule
        TEMPO$^{\bullet}$          & \good{$+228$} \\
        AQ semiquinone             & $+82$ \\
        viologen MV$^{+\bullet}$   & $+55$ \\
        \bottomrule
      \end{tabular}
      \end{table}
      \footnotesize TEMPO --- the textbook persistent radical (our validation anchor) ---
      ranks most stable, as it must. Exact from existing energies; no new calculation.
    \end{column}
    \begin{column}{0.46\textwidth}
      \textbf{Dimerization} $2\,\text{MV}^{+\bullet}\!\rightarrow[\text{MV}_2]^{2+}$
      (the real viologen failure mode):
      \begin{block}{$\Delta G_\text{dim} = +0.49$ eV \ (flagged)}
        Config sampling: minimum well-defined ($\sigma_\text{config}\!\approx$ 2 meV).
        \\ The uncertainty is the \bad{$+2$ pimer desolvation} --- the continuum's weak spot.
      \end{block}
      \vspace{-0.5ex}
      \footnotesize $\Rightarrow$ use dimerization as a \hi{relative} ranking; lean on
      disproportionation for \hi{ranking} (absolute values over-estimated, FINDINGS \#24).
    \end{column}
  \end{columns}
\end{frame}

% ---------------------------------------------------------------
\begin{frame}{Kinetics: inner- and outer-sphere reorganization}
  \begin{columns}[T]
    \begin{column}{0.44\textwidth}
      \scriptsize
      \textbf{Inner-sphere} $\lambda_i$ --- solute geometry relaxation, \textbf{4-point (Nelsen)}:
      \[\lambda_i=[E_O(q_R)\!-\!E_O(q_O)]+[E_R(q_O)\!-\!E_R(q_R)]\]
      gas single points; same four-point formula as D3TaLES. \good{Validated}.
      \vspace{0.5ex}

      \textbf{Outer-sphere} $\lambda_o$ --- solvent reorganization (Marcus/Born, \hi{Sharma Eq.\,13}):
      \[\lambda_o=\frac{z^2e^2}{8\pi\varepsilon_0 r_\text{hyd}}\!\left(\frac{1}{\varepsilon_\text{op}}-\frac{1}{\varepsilon_r}\right)\]
      MeCN ($\varepsilon_r{=}37.5,\ \varepsilon_\text{op}{=}1.81$); the production annotation
      uses a molecular-cavity nonequilibrium PCM, with Born as a check.
      \vspace{0.5ex}

      Report conventions explicitly: $\lambda_\text{het}=(\lambda_O+\lambda_R)/2+\lambda_{o,1}$;
      $\lambda_\text{SE}(d)=\lambda_O+\lambda_R+2\lambda_{o,1}(1-a/d)$.
    \end{column}
    \begin{column}{0.56\textwidth}
      \begin{center}
      \includegraphics[width=\linewidth]{d3tales_validation.png}
      \end{center}
      \vspace{-1.5ex}
      \centering\footnotesize
      At D3TaLES's own level, inner-sphere $\lambda$ agrees to
      \good{median $|\Delta|=0.09$ eV} ($n=322$).
    \end{column}
  \end{columns}
\end{frame}

% ---------------------------------------------------------------
\begin{frame}{Candidates: SA vs.\ reorganization energy}
  \begin{center}
  \includegraphics[height=0.55\textheight]{candidates_2x2_standalone_grafted.png}
  \end{center}
  \vspace{-2ex}
  \centering\scriptsize\textcolor{ink!70}{All six are \emph{grafted} model-compounds (redox core
  $+$ 4-methylbenzyl monomer proxy); short class names label them.}\par
  \vspace{0.4ex}
  \scriptsize
  \begin{columns}[T]
    \begin{column}{0.5\textwidth}
      \begin{itemize}\scriptsize
        \item one point per molecule (mean over its charge states); the six cluster tightly
              ($\lambda_i\!\approx$0.45--0.54 eV --- similar across families).
        \item \good{ammonium-NDI}: lowest $\lambda$ $\Rightarrow$ fastest ET; \hi{quinones} highest.
      \end{itemize}
    \end{column}
    \begin{column}{0.5\textwidth}
      \begin{itemize}\scriptsize
        \item gas $\to$ MeCN: $+\lambda_o$ lifts each point $\sim$0.5--0.6 eV (note per-panel y-zoom).
        \item $\lambda_i$ validated (D3TaLES $R^2{=}1.0$); $\lambda_o$ Marcus/Born estimate
              ($r_\text{hyd}$ caveat) --- read \emph{relative} trends.
      \end{itemize}
    \end{column}
  \end{columns}
\end{frame}

% ---------------------------------------------------------------
\begin{frame}{Functionalization $\times$ solvation: the 2$\times$2}
  \begin{columns}[T]
    \begin{column}{0.62\textwidth}
      \begin{center}
      \includegraphics[height=0.68\textheight,width=\linewidth,keepaspectratio]{candidates_2x2_standalone_grafted.png}
      \end{center}
    \end{column}
    \begin{column}{0.38\textwidth}
      \footnotesize
      \textbf{Rows} = \hi{before} (standalone core) vs \hi{after} functionalization
      (grafted: redox core $+$ 4-methylbenzyl monomer handle).\\
      \textbf{Cols} = gas ($\lambda_i$) vs MeCN ($\lambda_i{+}\lambda_o$).
      \vspace{1.2ex}

      \textbf{Functionalization barely moves $\lambda_i$} --- $\pm$15--50 meV, mixed sign:
      \begin{itemize}\scriptsize
        \item viologens: $+$15 / $+$29 meV
        \item PMDI $+$22; \good{NDI $-$44}; quinones $0$ / \good{$-$52}
      \end{itemize}
      \vspace{0.5ex}
      \textbf{Solvation} adds a near-uniform \hi{$+$0.5--0.6 eV} ($\lambda_o$).
      \vspace{1.2ex}

      $\Rightarrow$ the \good{redox core sets $\lambda$}; the monomer handle is a small
      perturbation, so the cheap standalone $\lambda$ is a good predictor of the grafted value.
    \end{column}
  \end{columns}
\end{frame}

% ---------------------------------------------------------------
\begin{frame}{What we learned (the durable lessons)}
  \begin{columns}[T]
    \begin{column}{0.5\textwidth}
      \textbf{\textcolor{good}{What works}}
      \begin{itemize}\footnotesize
        \item ranking (within charge class)
        \item xtb RRHO thermal (open-shell safe)
        \item disproportionation stability
        \item reorganization energy $\lambda$
        \item intact/bound state as a necessary-condition gate
        \item capacity ($n$, MW): exact
      \end{itemize}
    \end{column}
    \begin{column}{0.5\textwidth}
      \textbf{\textcolor{bad}{Negative results (saved effort)}}
      \begin{itemize}\footnotesize
        \item continuum ion-pairing: \bad{refuted} (spacing off by $\sim$11 V)
        \item explicit microsolvation (4 MeCN): \bad{no gain} for viologen
        \item absolute redox offset: \bad{unresolved} (OROP vs in-house anchors)
        \item dimerization absolute: \bad{unreliable} ($+2$ desolvation)
        \item solubility / crossover: proxies only
      \end{itemize}
    \end{column}
  \end{columns}
  \vspace{1ex}
  \begin{center}\small
  \textcolor{ink}{Corroborated by the metabolic-$\Delta G$ pipeline: 0.14 eV via isodesmic
  cancellation $\Rightarrow$ \hi{cancellation, not more physics, is the lever}.}
  \end{center}
\end{frame}

% ---------------------------------------------------------------
\begin{frame}{Putting it together: a tiered, multi-objective search}
  \begin{columns}[T]
    \begin{column}{0.46\textwidth}
      \textbf{Tiered funnel} (explicit solvation is a finisher, never for thousands):
      \begin{center}
      \begin{tikzpicture}[font=\scriptsize,node distance=3mm,
        t/.style={rectangle,rounded corners,draw=ink,fill=soft,thick,align=center,
                  minimum height=7mm}]
        \node[t,text width=52mm] (t0) {Tier 0 --- cheap rank (thousands)};
        \node[t,text width=38mm,below=of t0] (t1) {Tier 1 --- DFT+SMD (hundreds)};
        \node[t,text width=24mm,below=of t1,fill=teal!20] (t2) {Tier 2 --- explicit (dozens)};
        \draw[-{Stealth[length=2mm]},thick,ink] (t0)--(t1);
        \draw[-{Stealth[length=2mm]},thick,ink] (t1)--(t2);
      \end{tikzpicture}
      \end{center}
    \end{column}
    \begin{column}{0.54\textwidth}
      \textbf{Pareto scorecard --- built correctly}
      \begin{itemize}\footnotesize
        \item objectives = \good{trustworthy} axes only
              (capacity, $\lambda$, disprop.\ stability, ranked potential)
        \item proxies (solubility, SA, crossover) = \hi{filters}, not axes
        \item \textbf{integrity} = hard necessary-condition pre-filter
        \item missing applicable data = \textbf{INCOMPLETE}, never a favorable value
        \item split \textbf{anolyte / catholyte} (potential isn't ``more is better'')
        \item \textbf{$\sigma$-aware} domination (don't rank within error)
        \item read the front against the \emph{raw} axes --- no scalarized score (arbitrary
              weights, $\sigma$-blind, pool-relative)
      \end{itemize}
    \end{column}
  \end{columns}
\end{frame}

% ---------------------------------------------------------------
\begin{frame}{Applying it: candidate scorecard $\rightarrow$ shortlist}
  \footnotesize
  $\sigma$-aware Pareto on the current rankable grafted set; rows missing an applicable
  objective are \textbf{INCOMPLETE} and cannot enter or dominate the primary front:
  \begin{columns}[T]
    \begin{column}{0.61\textwidth}
      \begin{table}\scriptsize
      \renewcommand{\arraystretch}{1.1}
      \begin{tabular}{lrrrrc}
        \toprule
        id & E$_\text{an}$ & Cap & $\lambda$ & $\Delta G_\text{d}$ & Par. \\
        \midrule
        aq\_benzyloxy$^\star$         & $-1.46$ & 163 & 0.39 & 61 & \good{Y} \\
        pmdi                          & $-1.46$ & 160 & 0.47 & 64 & --- \\
        \good{mophquinone}$^\star$    & $-1.02$ & 176 & 0.52 & 92 & \good{Y} \\
        dmophquinone                  & $-1.00$ & 147 & 0.57 & 89 & --- \\
        \good{viologen}$^\dagger$     & $-0.94$ & 194 & 0.47 & 51 & TBA \\
        aq\_benzylamino$^\ddagger$    & $-1.57$ & 164 & ---  & 76 & partial \\
        \bottomrule
      \end{tabular}
      \end{table}
      \scriptsize Ordered by primary axis (E$_\text{an}$). E$_\text{an}$ [V vs Fc], Cap
      [mAh/g], $\lambda$ [eV], $\Delta G_\text{d}$ (disprop.) [kJ/mol].
    \end{column}
    \begin{column}{0.37\textwidth}
      \scriptsize $^\star$ = primary-front under both salts; $^\dagger$ joins under TBAPF$_6$;
      $^\ddagger$ is a partial front only (no QC-clean $\lambda$). The front and raw axes are
      reported without an arbitrary scalar score.
      \vspace{1ex}

      \hi{Completeness matters}: a missing or QC-flagged $\lambda$ cannot make a candidate look
      artificially fast. The row remains visible as a partial front, with its missing objective.
      \vspace{1ex}

      $\Rightarrow$ across this tight set voltage is \emph{non-discriminating}; \good{capacity
      $+$ disproportionation stability $+$ $\lambda$} drive the ranking.
    \end{column}
  \end{columns}
\end{frame}

% ---------------------------------------------------------------
\begin{frame}{Current focus: six multi-electron starting candidates}
  \begin{center}
  \includegraphics[height=0.78\textheight]{before_after_grafting.png}
  \end{center}
  \vspace{-1ex}
  \centering\scriptsize
  Two \hi{viologens}, two \hi{imides}, and two \hi{quinones}: capped core (left) versus
  Merrifield-grafted proxy (right).
\end{frame}

% ---------------------------------------------------------------
\begin{frame}{Next steps}
  \textbf{\textcolor{good}{Done}}: protocol-addressed energies; unified $\sigma$-tagged
  scorecard; integrity and completeness contracts; $\sigma$-aware Pareto engine.
  \vspace{1ex}
  \begin{enumerate}
    \item \hi{\textbf{Resolve the E$^\circ$ offset}} --- cross-run in-house and OROP molecules;
          audit OROP's thermodynamic cycle and reference convention before calibration.
    \item \textbf{Broaden benchmarks} --- second reductions and at least ten two-wave molecules.
    \item \textbf{Complete kinetics} --- conformer-matched core $\lambda$ for tether jumps and
          solvated $\lambda$ where a gas dianion is unbound.
    \item \textbf{Polymer effects} --- backbone/tether constraints \& inter-monomer
          dimerization (where the \emph{relative} dimerization ranking becomes key).
    \item \textbf{Absolute accuracy} (if needed) --- same-charge referencing to remove the
          per-class redox bias.
  \end{enumerate}
\end{frame}

% ---------------------------------------------------------------
\begin{frame}[plain]
  \vfill
  \begin{center}
    {\Large\textbf{\textcolor{ink}{Historical summary}}}\\[2ex]
    \begin{minipage}{0.85\textwidth}\centering
    We can \good{rank} redox candidates and screen \good{stability}, \good{kinetics}, and
    \good{state integrity} on a physics-based footing --- with each axis \hi{honestly labelled}
    by how far we can trust it.\\[2ex]
    The $\sigma$-aware scorecard \& Pareto engine are built and applied.
    Next: resolve absolute E$^\circ$ and close the incomplete kinetics rows before scaling.
    \\[1ex]\scriptsize Historical snapshot; current candidate and property validation remains separate.
    \end{minipage}
  \end{center}
  \vfill
\end{frame}

